\documentclass[11pt]{article}

\usepackage[T1]{fontenc}
\usepackage{lmodern}
\usepackage{microtype}
\usepackage{amsmath,amssymb,amsthm,mathtools}
\usepackage{booktabs}
\usepackage{tabularx}
\usepackage{enumitem}
\usepackage{xcolor}
\usepackage[margin=1.05in]{geometry}
\usepackage{fancyhdr}
\usepackage[hidelinks]{hyperref}

\definecolor{navy}{HTML}{17324D}
\definecolor{slate}{HTML}{435466}
\definecolor{paperblue}{HTML}{EAF2F8}
\definecolor{paperamber}{HTML}{FFF5D6}

\hypersetup{
  colorlinks=true,
  linkcolor=navy,
  citecolor=navy,
  urlcolor=navy,
  pdftitle={Marked circuits and fixed-five rooted interfaces in cubic graphs},
  pdfauthor={Atharva Vaidya}
}

\pagestyle{fancy}
\fancyhf{}
\fancyhead[L]{\small\color{slate}Marked circuits and rooted interfaces}
\fancyhead[R]{\small\color{slate}Research draft}
\fancyfoot[C]{\thepage}
\setlength{\headheight}{14pt}

\newtheorem{theorem}{Theorem}[section]
\newtheorem{proposition}[theorem]{Proposition}
\newtheorem{lemma}[theorem]{Lemma}
\newtheorem{corollary}[theorem]{Corollary}
\newtheorem{claim}[theorem]{Claim}
\theoremstyle{definition}
\newtheorem{definition}[theorem]{Definition}
\theoremstyle{remark}
\newtheorem{remark}[theorem]{Remark}

\newcommand{\cT}{\mathcal T}
\newcommand{\dotcupop}{\mathbin{\dot\cup}}

\title{\vspace{-1.4cm}\textbf{\color{navy}
Marked circuits and fixed-five rooted interfaces in cubic graphs}\\[0.35em]
\large Structural reductions and finite censuses}
\author{Atharva Vaidya}
\date{Research draft, 27 July 2026}

\begin{document}
\maketitle

\begin{abstract}
Let \(H\) be a finite simple cubic graph and let \(S\) be a
four-edge matching.  As algebraic setup, we first identify five
Eulerian edge sets covering every edge twice with an
everywhere-anisotropic flow in a four-dimensional quadratic space over
\(\mathbb F_2\), including the exact convention for loops.  We then
record six narrowly scoped results.
First, if \(H\) has a Tait colouring in which the four marks have one
common colour, and
\[
  |\delta_H(X)|+|S\cap E(H[X])|\geq 4
\]
whenever \(H[X]\) contains a cycle, then \(S\) is contained in an even
edge set whose every circuit component contains an even number of marks.
The proof combines elementary shore lemmas with three published
prescribed-cycle and decomposition results.  It does not require
universal separation after the common-colour Tait colouring has been
obtained.

Second, suppose \(H\) is cyclically \(4\)-edge-connected, the marks are
separated in a fixed common-colour Tait colouring, and a differently
coloured edge \(f\notin S\) is forbidden.  If no cycle of \(H-f\)
contains all four marks, then \(f\) belongs to an independent cyclic
four-edge cut.  Hence a cycle through all four marks and avoiding \(f\)
always exists when \(H\) is cyclically \(5\)-edge-connected.

Third, for the surviving exceptional rooted three-pole interface in a
separate \(D_5=\binom{[5]}2\) flow model, we prove a necessary
cycle-support obstruction.  If a nonbridge root has only one possible
label after the connector triangle is fixed, every cycle through that
root has a label support which spans all five coordinates and is either
non-bipartite or the star \(K_{1,4}\).  An exact gluing factorization
shows how this condition arises when a fixed-five boundary signature
equals one of the exceptional type patterns in the literature.  Missing
root base pairs also force explicit odd proper-core cuts.  These are
necessary restrictions, not eliminations.

Fourth, two exact solvers show that every nonempty nonbridge rooted
three-pole signature through order 17 contains one of three base pairs:
\(15\,645\,623\) roots in \(654\,676\) canonical cores are checked.
Fifth, on the 31 cyclically \(4\)-edge-connected bridgeless non-Tait
cubic graphs of order 22, all \(14\,322\) independent-edge deletion
poles have the full ten-type fixed-five boundary signature.

Sixth, an exhaustive canonical computation finds no
\emph{Tait-colourable} connected simple cubic graph on 22 vertices
having a universally separated four-edge matching.  Two separately
written enumeration methods agree on \(7\,174\,735\) colourable graphs
and \(95\,360\,112\) Tait colourings modulo global colour permutation.
This manuscript neither proves nor disproves the Five-Cycle Double Cover
Conjecture.
\end{abstract}

\noindent
\colorbox{paperamber}{%
\parbox{\dimexpr\textwidth-2\fboxsep\relax}{%
\textbf{AI-use disclosure.}
OpenAI Codex agents, under the direction of Atharva Vaidya, generated
the initial arguments, wrote the census programs, executed and audited
the computations, and drafted this manuscript.  Other Codex agents
independently rechecked selected proofs and enumerations; those are
AI-agent checks, not independent human verification or peer review.
All mathematical arguments used here are displayed for line-by-line
human checking, and the computational claims are separated from them.
The human author should independently verify the paper and establish
priority before submission.}}

\section{Introduction and exact scope}

A binary cycle of a graph is an edge set with even degree at every
vertex.  In a cubic graph its nonempty connected components are cycles.
The first part of this paper asks when four marked edges lie in a binary
cycle whose components each contain an even number of marks.  This
marked parity question arose in a larger investigation of the
Five-Cycle Double Cover Conjecture, but the results below do not resolve
that conjecture.

For clarity, the standard conjecture asks for at most five Eulerian
edge subsets whose incidence multiplicities cover each edge exactly
twice~\cite{Zhang2012,Oum2026}.  The orientable variant imposes additional directional
conditions and is a different problem.  Neither version is proved here.
In particular, our unrooted theorem has a marked-cut hypothesis, our
rooted theorem leaves a cyclic four-cut interface, and the census is
restricted to one order, simple graphs, and Tait-colourable inputs.
All \(D_5\) boundary signatures below use one fixed five-coordinate
universe.  They are not the unbounded-colour relation of
M\'a\v cajov\'a, Mazzuoccolo, and Tabarelli; our finite computations do
not classify colourings needing more than five coordinates.

\subsection*{The six contributions}

\begin{center}
\begin{tabularx}{\textwidth}{@{}
  >{\raggedright\arraybackslash\bfseries}p{0.18\textwidth}
  >{\raggedright\arraybackslash}X
  >{\raggedright\arraybackslash}p{0.25\textwidth}@{}}
\toprule
Result & Content & Nature of evidence\\
\midrule
Theorem~\ref{thm:core}
  & Four common-colour marks satisfying the marked cyclic-cut
    inequality admit an even-marked binary cycle.
  & Human proof using explicitly identified imported theorems.\\
Theorem~\ref{thm:gate}
  & A failed forbidden-edge reduction is localized to an independent
    cyclic four-cut; cyclic connectivity five closes the rooted case.
  & Human proof; only the four-independent-edge cycle theorem is
    imported.\\
Theorem~\ref{thm:translation-obstruction}
  & A nonbridge singleton root in the normalized \(D_5\) three-pole
    interface forces every root cycle to have a spanning translation
    blocker.
  & Human algebraic and graph-cut proof; a necessary obstruction only.\\
Theorem~\ref{thm:base-pair-certificate}
  & Failure of three rooted base pairs forces an odd proper-core cut
    carrying one common coordinate.
  & Elementary human proof; a necessary certificate only.\\
Theorem~\ref{thm:order17-base-pair}
  & Every nonempty nonbridge rooted three-pole signature through order
    17 contains a base pair.
  & Two exact solvers with matching per-shard transcript digests.\\
Theorem~\ref{thm:cyclic4-cap-census}
  & Every order-22 independent-edge deletion in the cyclically
    four-connected non-Tait cap corpus has the full fixed-five signature.
  & Two full boundary classifiers with byte-identical transcripts.\\
Theorem~\ref{thm:census}
  & No Tait-colourable connected simple cubic graph of order 22 has a
    universally separated four-edge matching.
  & Exhaustive computation with two algorithmically different
    enumerators.\\
\bottomrule
\end{tabularx}
\end{center}

The human arguments appear to be useful deductions from known cycle and
boundary-signature machinery.  The rooted factorizations and finite
frontiers appear new to the authors, but the priority search performed
for this draft was targeted rather than exhaustive.

\subsection*{A scope correction discovered during audit}

An earlier project note stated the computational result over all
connected simple cubic graphs of order 22.  The programs, correctly for
the motivating application, search only graphs having at least one Tait
colouring.  This qualifier is logically essential.  If a graph has no
Tait colouring, the assertion ``separated in every Tait colouring'' is
vacuously true under the literal universal quantifier.  Moreover, every
cubic graph on 22 vertices has a matching of size at least four: a
maximal matching of size \(m\) has endpoints dominating at most \(6m\)
vertices, so \(m\geq\lceil22/6\rceil=4\).  Theorem~\ref{thm:census}
therefore states exactly the nonvacuous result established by the
programs.

\section{Definitions and imported results}

All graphs are finite.  Unless stated otherwise, they are simple.  For
\(X\subseteq V(G)\), let \(\delta_G(X)\) be the set of edges with exactly
one end in \(X\), and let \(G[X]\) be the induced subgraph.  A
\emph{cycle} is a connected 2-regular subgraph.  A \emph{binary cycle}
is an edge subset \(Q\subseteq E(G)\) such that every vertex has even
degree in \(Q\).

A graph is \emph{cyclically \(k\)-edge-connected} if it has no edge cut
of size less than \(k\) with a cycle on both shores.  A
\emph{Tait colouring} of a cubic graph is a proper edge-colouring with
three colours \(a,b,c\).  Every two colour classes induce a spanning
2-regular subgraph, called a \emph{bichromatic factor}; its components
are even cycles.

\subsection*{An exact elliptic-flow reformulation}

For later use, put
\[
 \Gamma=\left\{x\in\mathbb F_2^5:\sum_{i=0}^4x_i=0\right\},
 \qquad
 q(x)=\sum_{0\leq i<j\leq4}x_ix_j.                    \tag{2.1}
 \label{eq:quadratic-space}
\]
If \(x\in\Gamma\), then
\[
 q(x)=\binom{\operatorname{wt}(x)}2\pmod2.
\]
Thus \(q(x)=1\) exactly for the ten vectors of weight two, which we
identify with
\[
 D_5=\binom{\{0,1,2,3,4\}}2.                          \tag{2.2}
 \label{eq:d5}
\]
The polar form is
\[
 B(x,y)=q(x+y)+q(x)+q(y)=\sum_{i=0}^4x_iy_i
 \qquad(x,y\in\Gamma).                                \tag{2.3}
\]
It is nondegenerate: a vector orthogonal to every \(e_i+e_j\) has all
coordinates equal, and the all-one vector is not in \(\Gamma\).
The five nonzero singular vectors are
\(\mathbf1+e_i\), \(0\leq i<5\).  They span \(\Gamma\), with their sum
as their only dependence.  Hence every isometry permutes them
faithfully.  Coordinate permutations already realize all five
permutations, so
\[
 O(\Gamma,q)\cong S_5.                                 \tag{2.4}
\]

\begin{proposition}[Elliptic-flow equivalence]
\label{prop:elliptic-flow}
Let \(G\) be a finite multigraph.  Five Eulerian edge subsets
\(C_0,\ldots,C_4\) cover each edge exactly twice if and only if there is
a map \(\phi:E(G)\to\Gamma\) such that
\[
 q(\phi(e))=1\quad(e\in E(G)),\qquad
 \sum_{\substack{\text{edge-ends}\\\text{incident with }v}}\phi(e)=0
 \quad(v\in V(G)).                                     \tag{2.5}
 \label{eq:anisotropic-flow}
\]
A loop is counted by its two edge-ends in the second sum.
\end{proposition}

\begin{proof}
Given the five edge subsets, let coordinate \(i\) of \(\phi(e)\) be
one exactly when \(e\in C_i\).  Exact double coverage makes
\(\phi(e)\) a weight-two vector, hence an element of \(q^{-1}(1)\).
Coordinate \(i\) of the vertex equation is precisely the degree in
\(C_i\) modulo two.  Conversely, \(q(\phi(e))=1\) forces
\(\phi(e)\) to have exactly two nonzero coordinates.  Defining
\(C_i=\{e:\phi(e)_i=1\}\) therefore gives exact double coverage, and
the five coordinate equations give the five even-degree conditions.
The two constructions are inverse.
\end{proof}

Parallel edges are distinct edge variables.  A loop contributes twice
to each selected degree, so it cancels in the characteristic-two
vertex equation; counting a loop only once in a SAT/XOR incidence row
would be an incorrect encoding.  Empty \(C_i\) are allowed, so this
five-coordinate formulation means ``at most five'' in the standard
sense~\cite{Zhang2012,Oum2026}.  It imposes no orientation condition.
The underlying two-subset encoding also appears as a
\emph{CDC-colouring} in M\'a\v cajov\'a, Mazzuoccolo, and
Tabarelli~\cite[Definition 3.1]{MacajovaEtAl2026}.
Proposition~\ref{prop:elliptic-flow} is an exact elementary
reformulation, not a resolution or a priority claim; priority for the
quadratic packaging has not been established.

\begin{definition}[Separation]
Let \(S\) be a matching in a Tait-coloured cubic graph.  The marks are
\emph{separated in that colouring} if every bichromatic factor cycle
contains at most one member of \(S\).  If the graph is Tait-colourable,
then \(S\) is \emph{universally separated} if it is separated in every
Tait colouring.
\end{definition}

We use the following three published results.  They are the only
non-elementary imported ingredients in the proof of
Theorem~\ref{thm:core}.

\begin{enumerate}[label=(I\arabic*),leftmargin=2.8em]
\item\label{imp:kp}
  Knappe and Pitz~\cite[Fact 5.4]{KnappePitz2020} prove \(g(3)=3\).
  In the form used here: if \(G\) is 3-edge-connected and
  \(F\subseteq E(G)\) has at most three edges, then a circuit contains
  \(F\) if and only if \(F\) contains no odd edge cut of \(G\).
\item\label{imp:aeeh}
  Aldred, Ellingham, Hemminger, and Holton
  \cite[Theorem 3.3]{AldredEtAl1997} prove that a free edge system of
  size at most four in a quasi \(4\)-connected graph lies on a cycle.
  In particular, four independent edges lie on a cycle.  They also
  note that cyclically \(4\)-edge-connected cubic graphs are quasi
  \(4\)-connected.
\item\label{imp:decomp}
  Every 3-connected cubic graph is obtained by repeated vertex
  3-sums of cyclically \(4\)-edge-connected cubic graphs.  This is
  Theorem 6.1 of Nedela, Seifrtova, and Skoviera
  \cite{NedelaEtAl2024}; see also Pastor's component-tree theorem
  \cite{Pastor2018}.
\end{enumerate}

Here a circuit in \ref{imp:kp} is a connected closed trail.  In a cubic
graph its edge-induced degrees are positive and even, hence equal to
two, so its edge set is a cycle.

\begin{lemma}[Tait cut parity]\label{lem:parity}
In a Tait-coloured cubic graph, the numbers of cut edges of colours
\(a,b,c\) have the same parity.  In particular, the three edges of a
three-edge cut have three distinct colours, and a Tait-colourable cubic
graph has no bridge.
\end{lemma}

\begin{proof}
For a colour \(p\), counting incidences of the \(p\)-coloured perfect
matching with \(X\) gives
\[
  |X|\equiv |\delta_G(X)\cap E_p|\pmod 2.
\]
This holds for all three colours.  If the cut has size three, all three
counts are odd and sum to three, so each is one.  If a single edge were
a cut, its colour count would have parity different from the other two.
\end{proof}

\begin{lemma}[Precolouring separated marks]\label{lem:precolour}
Let \(H\) be Tait-colourable and let \(S\) be universally separated.
Every map \(S\to\{a,b,c\}\) extends to a Tait colouring.  In particular,
there is a Tait colouring in which all marks have one common colour.
\end{lemma}

\begin{proof}
Start with any Tait colouring.  If a mark \(e\) currently has colour
\(p\) and is requested to have colour \(q\), interchange \(p\) and
\(q\) on the unique \(pq\)-factor cycle through \(e\).  Universal
separation says that this cycle contains no other mark, so all previously
fixed marks keep their colours.  Process the marks one at a time.
\end{proof}

\section{Two elementary shore lemmas}

The marked hypothesis used throughout the first theorem is
\begin{equation}
 |\delta_H(X)|+|S\cap E(H[X])|\geq4
 \quad\text{whenever \(H[X]\) contains a cycle.}
 \tag{MC}\label{eq:mc}
\end{equation}

\begin{lemma}[Two-mark shore]\label{lem:two-shore}
Let \(A\subsetneq V(H)\), put \(P=H[A]\), and suppose that \(P\) is
connected and subcubic with minimum degree two.  Suppose
\(|\delta_H(A)|\in\{2,3\}\), the boundary edges have distinct ends in
\(A\), and exactly two marks lie in \(E(P)\).  If
\eqref{eq:mc} holds for every cyclic \(X\subseteq A\), then one cycle of
\(P\) contains both marks.
\end{lemma}

\begin{proof}
Delete the bridges of \(P\) and form the tree of bridge-components.  If
\(P\) has no bridge, it is 2-edge-connected.  A 2-edge-connected
subcubic graph of minimum degree two has no cutvertex: two components
behind a cutvertex would each require at least two incident edges at
that vertex, giving degree at least four.  Thus \(P\) is 2-connected.
Subdivide the two marked edges.  Any two vertices of a 2-connected graph
lie on a common cycle, and suppressing the subdivisions gives a cycle
through both marks.

Now suppose the bridge-component tree is nontrivial.  Every leaf
component \(K\) is cyclic; a single isolated vertex would have only one
incident bridge but degree at least two in \(P\).  Let \(p(K)\) count
the pole-boundary edges incident with \(K\), and let \(s(K)\) count its
internal marks.  Exactly one bridge leaves \(K\), so \eqref{eq:mc} gives
\[
  1+p(K)+s(K)\geq4.
\]
If \(K\) does not contain both marks, then \(s(K)\leq1\), hence
\(p(K)\geq2\).  Two such leaves would consume at least four distinct
pole-boundary edges, impossible because there are at most three.
Therefore a leaf contains both marks.  That bridge-component is
2-edge-connected, subcubic, and of minimum degree two, so the preceding
argument supplies the desired cycle.
\end{proof}

\begin{lemma}[One-mark shore]\label{lem:one-shore}
Under the setup of Lemma~\ref{lem:two-shore}, suppose instead that
\(|\delta_H(A)|=3\), its ends in \(A\) are \(u_1,u_2,u_3\), and exactly
one mark \(e\) lies in \(E(P)\).  If \eqref{eq:mc} holds on every cyclic
\(X\subseteq A\), then for every pair \(i\neq j\) there is a
\(u_i\)-\(u_j\) path in \(P\) containing \(e\).
\end{lemma}

\begin{proof}
If \(P\) had a bridge, a leaf \(K\) of its bridge-component tree would
be cyclic.  With \(p(K)\) and \(s(K)\in\{0,1\}\) as above,
\[
  1+p(K)+s(K)\geq4.
\]
An unmarked leaf consumes all three pole-boundary edges, while a marked
leaf consumes at least two.  Two leaves cannot do so using only three
boundary edges.  Hence \(P\) is 2-edge-connected and, by the subcubic
cutvertex argument, 2-connected.

Fix \(i\neq j\), add an auxiliary edge \(u_i u_j\), and use the standard
fact that any two edges of a 2-connected graph lie on a common cycle.
The assertion remains valid if the auxiliary edge is parallel to an
existing edge.  A cycle containing the auxiliary edge and \(e\), with
the auxiliary edge deleted, is the required path.
\end{proof}

\begin{corollary}[Low-cut reductions]\label{cor:lowcut}
If a graph satisfying \eqref{eq:mc} has no binary cycle containing all
four marks with an even number of marks in each component, then it has
no cyclic two-edge cut.  Moreover, an unmarked cyclic three-edge cut
cannot split the marks \(2+2\).
\end{corollary}

\begin{proof}
For a cyclic two-cut, \eqref{eq:mc} puts at least two internal marks on
each shore.  Since there are four marks in total, the cut is unmarked
and each shore has exactly two.  Lemma~\ref{lem:two-shore} gives one
two-mark cycle on each shore, whose union is a valid binary cycle.  The
same argument applies to the two shores of an unmarked \(2+2\) cyclic
three-cut.
\end{proof}

\section{The four-mark core theorem}

\begin{theorem}[Four-mark core closure]\label{thm:core}
Let \(H\) be a connected simple cubic graph and let \(S\) be a
four-edge matching.  Assume:
\begin{enumerate}[label=(\roman*),leftmargin=2.4em]
\item \(H\) has a Tait colouring in which every member of \(S\) has one
      common colour \(c\); and
\item the marked cyclic-cut inequality \eqref{eq:mc} holds.
\end{enumerate}
Then there is a binary cycle \(Q\subseteq E(H)\) containing \(S\) such
that every nonempty component of \(Q\) contains an even number of marks.
\end{theorem}

\begin{remark}
Universal separation is not a hypothesis of
Theorem~\ref{thm:core}.  By Lemma~\ref{lem:precolour}, it is one
sufficient way to obtain hypothesis (i).  The proof below uses only the
one fixed common-colour Tait colouring.
\end{remark}

\begin{proof}
Suppose that no such \(Q\) exists.  By
Corollary~\ref{cor:lowcut}, \(H\) has no cyclic two-edge cut.

\smallskip
\noindent\emph{Step 1: \(H\) is 3-connected.}
Lemma~\ref{lem:parity} excludes bridges.  A cutvertex in a bridgeless
cubic graph is impossible, because one component behind the cutvertex
would receive only one of its three incident edges, making that edge a
bridge.

Suppose \(\{u,v\}\) is a two-vertex cut.  Every component \(D\) of
\(H-\{u,v\}\) has at least three boundary edges; otherwise the cubic
degree sum makes both shores of \(\delta_H(D)\) cyclic, contrary to the
absence of a cyclic cut of size at most two.  There are at least two
components and at most six boundary incidences at \(u,v\).  Hence
\(uv\notin E(H)\), there are exactly two components \(D_1,D_2\), and
each has three boundary edges.  Neither component sends all three edges
to only one of \(u,v\), since the other would then be a cutvertex.
Relabel so that \(D_1\) sends two edges to \(u\) and one to \(v\).
Then
\(\delta_H(D_1\cup\{u\})\) has size two.  Both shores are cyclic by the
cubic degree count and simplicity, a contradiction.  Thus \(H\) is
3-connected.

\smallskip
\noindent\emph{Step 2: eliminate marked cyclic three-cuts.}
Let \(B=\delta_H(A)\) be a cyclic three-edge cut and write
\[
 r=|B\cap S|,\quad
 a=|S\cap E(H[A])|,\quad
 d=|S\cap E(H[V(H)\setminus A])|.
\]
The inequality \eqref{eq:mc} on the two shores gives \(a,d\geq1\), and
\[
 a+d+r=4. \tag{4.1}\label{eq:marks-split}
\]
In particular, \(r\leq2\).  Since a shared boundary end would expose a
cyclic cut of size at most two, the three cut edges have distinct ends
on each shore.

If \(r=2\), then \(a=d=1\).  On each shore,
Lemma~\ref{lem:one-shore} gives a path through its unique internal mark
between the ends of the two marked cut edges.  The two paths and those
two cut edges form one cycle containing all four marks.

Suppose \(r=1\).  Up to exchanging shores, \(a=1\) and \(d=2\).  Cap
the two-mark shore by adjoining a new cubic vertex \(z\) at its three
boundary ends; call the resulting cubic graph \(\widehat D\).
Lemma~\ref{lem:parity} shows that the fixed Tait colouring caps
properly.  The cap is 3-edge-connected.  Indeed, it has no bridge, and
a two-edge cut in a cubic graph has two cyclic shores; the shore
avoiding \(z\) would induce a cyclic two-cut of \(H\).

Let \(F\) consist of the cap edge corresponding to the marked member of
\(B\) and the two internal marks.  These three edges are independent.
They contain no odd cut.  A one-edge cut is absent.  If the three edges
formed a three-edge cut, that cut could not be a vertex star because
\(F\) is independent.  For a nontrivial three-cut, take the cyclic
shore avoiding \(z\).  Both internal marks would be boundary edges and
there would be no internal mark, contradicting \(3+0\geq4\) in
\eqref{eq:mc}.  Imported result \ref{imp:kp} therefore gives a cycle
through \(F\).

At \(z\), that cycle uses the marked cap edge and one other cap edge.
Deleting \(z\) leaves a path through both internal marks between the
corresponding boundary positions.  On the one-mark shore,
Lemma~\ref{lem:one-shore} supplies a path through its mark for the same
boundary pair.  Gluing with the two selected cut edges gives one
four-mark cycle.  Both cases contradict the assumed failure.  Hence
every cyclic three-cut that remains is unmarked.

By Corollary~\ref{cor:lowcut}, an unmarked cyclic three-cut cannot split
the marks \(2+2\).  Equation~\eqref{eq:marks-split} shows that every
such cut splits them \(1+3\).

\smallskip
\noindent\emph{Step 3: organize all remaining cuts at once.}
Apply imported result \ref{imp:decomp}.  Split along cyclic three-cuts,
cap each shore by a root vertex, and continue until the factors are
cyclically \(4\)-edge-connected.  Reverse reconstruction gives a tree
\(\mathcal T\): its nodes are the final factors and its edges record
the vertex 3-sums.  Each tree edge corresponds to a principal cyclic
three-cut of \(H\).

All principal cuts are unmarked, so every marked edge belongs to one
factor.  Give a node weight equal to its number of marks.  Deleting any
tree edge splits the total weight four as \(1+3\).  There is a node
\(K\) for which every component of \(\mathcal T-K\) has weight one.
To prove this, begin at any node and move into a component of weight
three whenever one exists.  The walk never backtracks and must stop.
At the stopping node no component has weight three; weights two and
zero are excluded respectively by the \(1+3\) rule and by positivity
on both sides of every tree edge.

Let \(K\) contain \(r\) marked edges.  It has
\(q=4-r\) incident one-mark branches, represented in the capped factor
by root vertices \(z_1,\ldots,z_q\).  Transfer the fixed Tait colouring
to the cap.  The \(r\) central marks have colour \(c\); at each \(z_j\),
let \(f_j\) be the unique incident edge of colour \(c\).  Let
\[
 F=(S\cap E(K))\cup\{f_1,\ldots,f_q\},
\]
interpreted as a set, so a shared edge between adjacent cap roots is
included once.  The set \(F\) is a matching, \(1\leq |F|\leq4\), it
contains every central mark, and an edge of \(F\) is incident with each
\(z_j\).

If \(|F|=4\), imported result \ref{imp:aeeh} supplies a cycle through
\(F\).  If \(|F|\leq3\), then \(F\) contains no odd cut: a one-cut is
excluded by 3-edge-connectivity, a trivial three-cut is not a matching,
and a nontrivial three-cut is excluded by cyclic
\(4\)-edge-connectivity.  Imported result \ref{imp:kp} again supplies a
cycle through \(F\).  Denote the resulting central cycle by \(Z\).

For every \(j\), the cycle \(Z\) passes through \(z_j\) and selects two
of its three boundary positions.  The corresponding entire branch of
\(\mathcal T-K\) is a shore with exactly one mark.
Lemma~\ref{lem:one-shore} gives a path through that mark for the selected
boundary pair.  Substitute these paths while reversing the 3-sums.
Each substitution preserves a single cycle.  If two adjacent cap roots
share the selected edge \(z_i z_j\), reversing both sums turns that
single edge into the single boundary edge between the two expanded
shores; it is not duplicated.  The final cycle contains the \(r\)
central marks and one mark from each of the \(q\) branches, hence all
four marks.  This final contradiction proves the theorem.
\end{proof}

\begin{remark}[Why the conclusion allows two components]
An unmarked cyclic three-cut with a \(2+2\) mark split can be closed by
one two-mark cycle on each shore.  Nothing in the hypotheses forces
those two cycles to merge.  Thus the binary-cycle conclusion is the
proved statement; an unconditional one-cycle strengthening is not
claimed.
\end{remark}

\section{A forbidden root edge and the four-cut gate}

We now use a stronger fixed-colouring hypothesis and obtain a rooted
conclusion in a more highly connected class.

\begin{theorem}[Independent four-cut gate]\label{thm:gate}
Let \(H\) be a finite simple cyclically \(4\)-edge-connected cubic graph.
Let \(S=\{s_1,s_2,s_3,s_4\}\) be a matching and let \(f\notin S\).
Assume that \(H\) has a Tait colouring in which
\begin{enumerate}[label=(\roman*),leftmargin=2.4em]
\item all members of \(S\) have one common colour \(c\);
\item \(f\) has a different colour; and
\item \(S\) is separated in this fixed colouring.
\end{enumerate}
If no cycle of \(H-f\) contains all four members of \(S\), then \(f\)
belongs to an independent cyclic four-edge cut of \(H\).
\end{theorem}

\begin{proof}
Write \(f=uv\),
\[
 N_H(u)-\{v\}=\{u_1,u_2\},\qquad
 N_H(v)-\{u\}=\{v_1,v_2\},
\]
and form the root-edge reduction
\[
 R=H\mathbin{\ominus}f
   :=H-\{u,v\}+\{\,\bar u=u_1u_2,\ \bar v=v_1v_2\,\}.
 \tag{5.1}\label{eq:reduction}
\]
The graph \(H\) has at least eight vertices because \(S\) is a
four-edge matching.  It has no triangle: the triangle boundary has size
three and its complementary shore is cyclic by the cubic degree count,
contrary to cyclic \(4\)-edge-connectivity.  Consequently the four
neighbours appearing in \eqref{eq:reduction} do not create loops,
parallel edges, or coincident new edge objects.  Also \(H\) is
3-connected (for cubic graphs, cyclic connectivity and ordinary
connectivity agree through order three), so deleting the adjacent pair
\(u,v\) leaves a connected graph.  Thus \(R\) is connected, simple, and
cubic.

Map every mark \(s\in S\) to an edge \(\bar s\in E(R)\) by
\[
 \bar s=
 \begin{cases}
  \bar u,&\text{if \(s\) is incident with \(u\)},\\
  \bar v,&\text{if \(s\) is incident with \(v\)},\\
  s,&\text{otherwise}.
 \end{cases}
 \tag{5.2}\label{eq:mark-image}
\]
The four images are distinct.  We claim that they form a matching.
Only adjacency between a new image and an unchanged image needs proof.
Suppose, without loss of generality, that \(s=uu_1\) is a mark of
colour \(c\), while \(f\) has colour \(a\).  Then \(uu_2\) has the
third colour \(b\).  An unchanged mark incident with \(u_2\) would lie
with \(s\) on the \(bc\)-factor cycle, contradicting fixed-colouring
separation.  A mark incident with \(u_1\) would already be adjacent to
\(s\).  The other cases are symmetric.

If \(R\) were cyclically \(4\)-edge-connected, imported result
\ref{imp:aeeh} would give a cycle through the four independent images.
Replace \(\bar u\), when used, by \(u_1uu_2\), and similarly replace
\(\bar v\).  The lifted cycle lies in \(H-f\) and contains every
original mark, contrary to hypothesis.  Hence \(R\) has a cyclic cut
\(\delta_R(X)\) of size \(k\leq3\).

We pull this cut back to \(H\).  If \(u_1,u_2\) lie on opposite shores,
then \(\bar u\) is a cut edge; placing \(u\) on either shore replaces
\(\bar u\) by exactly one crossing edge.  If \(u_1,u_2\) lie on one
shore, place \(u\) there and create no crossing edge.  The same choices
apply to \(v\).  Unless both neighbour pairs are unsplit and lie on
opposite shores, choose \(u,v\) on the same shore.  Then \(f\) is
internal, the cut size remains \(k\leq3\), and both shore cycles lift:
an internal new edge is replaced by its two-edge path, while a crossing
new edge was not part of either shore cycle.  This contradicts cyclic
\(4\)-edge-connectivity of \(H\).

Therefore \(u_1,u_2\) lie together on one shore and \(v_1,v_2\) lie
together on the other.  Restoring \(u,v\) on those forced shores gives
\[
 \delta_H(X^\uparrow)=\delta_R(X)\dotcupop\{f\}.
 \tag{5.3}\label{eq:pullback}
\]
Both shores remain cyclic, so \(k+1\geq4\).  Since \(k\leq3\), we have
\(k=3\).

It remains to prove independence.  If two edges of the cyclic
three-cut \(\delta_R(X)\) shared an endpoint \(w\) on one shore, then
only one edge at \(w\) would remain internal.  No cycle of that shore
uses \(w\).  Moving \(w\) to the other shore replaces two crossing
edges by one and preserves cycles on both shores, producing a cyclic
two-cut of \(R\).  The pullback argument just given excludes such a
cut, since it would produce a cyclic cut of size at most three in
\(H\).  Thus the three edges of \(\delta_R(X)\) are independent.  In
the forced configuration they avoid the new edges and hence avoid
\(u,v\).  Together with \(f\), they form an independent cyclic
four-edge cut of \(H\).
\end{proof}

\begin{corollary}[Rooted closure at cyclic connectivity five]
\label{cor:cyclic5}
Under the hypotheses of Theorem~\ref{thm:gate}, if \(H\) is cyclically
\(5\)-edge-connected, then \(H-f\) has a cycle containing all four
marks.
\end{corollary}

\begin{proof}
Otherwise Theorem~\ref{thm:gate} produces a cyclic four-edge cut.
\end{proof}

\begin{remark}[What remains at cyclic connectivity four]
Theorem~\ref{thm:gate} is a one-way obstruction localization, not a
characterization: a root-containing independent four-cut need not
prevent an avoiding cycle.  The unresolved case is the corresponding
four-pole interface.  No unrestricted rooted theorem follows from the
gate alone.
\end{remark}

\section{Cycle translations in an exceptional rooted interface}

This section records a second, algebraic restriction on a rooted
interface.  Its ``root'' is a distinguished cap edge whose possible
\(D_5\)-labels are studied; it is not the forbidden edge of
Theorem~\ref{thm:gate}.  The conclusion is a necessary obstruction, not
a proof that the exceptional interface is impossible.  Throughout this
section, \(\mathcal R_5\) and \(\mathcal C_5\) denote exact signatures
inside one fixed five-coordinate \(D_5\) universe.  This is narrower
than the published unbounded-colour relation.

\subsection{The \texorpdfstring{\(D_5\)}{D5} model}

Recall \(\Gamma\), \(q\), and \(D_5=q^{-1}(1)\) from
\eqref{eq:quadratic-space}--\eqref{eq:d5}.  We identify a member of
\(D_5\) with its weight-two incidence vector and write addition in
\(\Gamma\) as symmetric difference.  A
\(D_5\)-labelling of a cubic multipole assigns a member of \(D_5\) to
every proper edge and semiedge so that the three incident labels sum to
zero at each internal vertex.  Such a local triple is exactly the
three-edge set of a triangle in \(K_5\).

Let \(Q\) be a cubic three-pole whose ordered connector word has been
normalized to
\[
 \tau=(01,02,12),                                      \tag{6.1}
 \label{eq:tau}
\]
and let \(r\) be a distinguished proper edge.  Its exact root signature
is
\[
 \mathcal R_5(Q,r)=
 \{A\in D_5:\text{\(Q\) has a \(D_5\)-labelling with word
 \(\tau\) and root label \(A\)}\}.                       \tag{6.2}
\]
The stabilizer of the ordered word \(\tau\) fixes \(0,1,2\) and may
swap \(3,4\).  Therefore \(\mathcal R_5(Q,r)\) is a union of the seven
orbits
\[
 \{01\},\ \{02\},\ \{12\},\
 \{03,04\},\ \{13,14\},\ \{23,24\},\ \{34\}.           \tag{6.3}
\]

\begin{lemma}[Cycle translation]\label{lem:cycle-translation}
Let \(\lambda\) be a normalized root labelling of \(Q\), let \(C\) be
a cycle of proper edges, and let \(h\in\Gamma\).  If
\[
 \lambda(e)+h\in D_5\qquad(e\in E(C)),                  \tag{6.4}
 \label{eq:translation-condition}
\]
then adding \(h\) to every label on \(C\), and changing no other label,
gives another normalized root labelling.
\end{lemma}

\begin{proof}
Every vertex of \(C\) is incident with exactly two translated edges, so
the change in its incident-label sum is \(h+h=0\).  All other vertices
are unchanged.  A proper cycle contains no connector semiedge, so the
word \(\tau\) in \eqref{eq:tau} remains fixed.  Condition
\eqref{eq:translation-condition} says that all translated values remain
in the allowed label set.
\end{proof}

\begin{lemma}[Translation blockers]\label{lem:blockers}
Let \(L\subseteq E(K_5)=D_5\) be nonempty, with its edges connected
under edge adjacency.  The only
\(h\in\Gamma\) for which
\[
 A+h\in D_5\qquad(A\in L)                              \tag{6.5}
 \label{eq:blocker-translation}
\]
is \(h=0\) if and only if the graph
\((\{0,1,2,3,4\},L)\)
\begin{enumerate}[label=(\roman*),leftmargin=2.4em]
\item spans all five coordinate vertices; and
\item is either non-bipartite or is the four-edge star \(K_{1,4}\).
\end{enumerate}
\end{lemma}

\begin{proof}
Every nonzero element of \(\Gamma\) has weight two or four.  First let
\(h\) have weight four, say
\(h=\{0,1,2,3,4\}\setminus\{v\}\).  For \(A\in D_5\), the symmetric
difference \(A+h\) has weight two exactly when \(A\) avoids \(v\).
Thus all weight-four translations are blocked exactly when \(L\)
spans all five vertices.

Now let \(h=B\) have weight two.  For \(A\in D_5\),
\[
 |A\mathbin\triangle B|=2
 \quad\Longleftrightarrow\quad |A\cap B|=1.
\]
Consequently \(B\) translates every member of \(L\) within \(D_5\)
exactly when every edge of \(L\) crosses the cut
\[
 B\mid(\{0,1,2,3,4\}\setminus B).                       \tag{6.6}
\]
If \(L\) spans and is non-bipartite, it is contained in no such cut.
If \(L\) is spanning and bipartite, edge-adjacency connectedness makes
it connected, so its bipartition is
unique.  A \(2+3\) bipartition gives a weight-two translation.  A
\(1+4\) bipartition has no side of size two; connectedness and spanning
force all four incident edges, namely \(K_{1,4}\).
The weight-two and weight-four cases prove the equivalence.
\end{proof}

\begin{remark}
Up to coordinate permutation, the inclusion-minimal connected blockers
are \(K_{1,4}\), the bull graph, and \(C_5\).  This follows directly:
a minimal spanning non-bipartite graph contains a triangle or a
five-cycle.  Minimality gives \(C_5\) in the latter case.  Starting
from a triangle, the unused vertices must attach at distinct triangle
vertices: attaching both to one vertex contains a star blocker, while a
pendant path leaves a triangle plus a disjoint edge after its joining
edge is deleted.  The remaining case is the bull.
\end{remark}

\begin{theorem}[Cycle-support obstruction]
\label{thm:translation-obstruction}
Suppose
\[
 \mathcal R_5(Q,r)=\{A\},\qquad
 A\in\{01,02,12,34\},
\]
and \(r\) is a nonbridge of the proper graph of \(Q\).  In every
normalized root labelling and on every proper cycle \(C\) containing
\(r\), the set of distinct labels
\[
 L_C=\{\lambda(e):e\in E(C)\}
\]
spans all five coordinates and is either non-bipartite or \(K_{1,4}\).
\end{theorem}

\begin{proof}
Consecutive cycle labels share one coordinate because they are two
edges of the local label triangle.  Thus \(L_C\) is nonempty and
connected under edge adjacency.  If it were not a blocker,
Lemma~\ref{lem:blockers} would give a nonzero \(h\) satisfying
\eqref{eq:translation-condition}.  By
Lemma~\ref{lem:cycle-translation}, translation on \(C\) gives another
normalized labelling whose root label is \(A+h\in D_5\).  Since
\(h\neq0\), this differs from \(A\), contradicting the singleton
signature.  Lemma~\ref{lem:blockers} gives the stated classification.
\end{proof}

\subsection{Exact factorization across the cyclic three-cut}

We now state the precise setting in which the singleton arises.  Let
\(G=P+\{e,f\}\) be a bridgeless cubic cap with no cyclic two-edge cut.
Assume that a cyclic three-edge cut, disjoint from \(e,f\), has connected
shores, splits the four terminals of \(P\) as \(2+2\), and contains one
whole cap edge on each shore.  Cutting the three cut edges gives rooted
three-poles \((Q_1,e)\) and \((Q_2,f)\).  Order corresponding connectors,
normalize them to \(\tau\), and write
\[
 R=\mathcal R_5(Q_1,e),\qquad S=\mathcal R_5(Q_2,f).
\]

For nonempty \(R,S\subseteq D_5\), let
\[
 \operatorname{rel}(R,S)\subseteq
 \{\mathsf E,\mathsf I,\mathsf D\}
\]
record which of equality, one-point intersection, and disjointness
occur among cross-pairs \((A,B)\in R\times S\).  Let \(\pi\) denote the
pairing of the four boundary semiedges induced by the deleted cap
edges.  In the standard four-pole notation, a doubled boundary word
\((A,A,B,B)\) has type
\[
 AA,\quad AT_\pi,\quad\text{or}\quad T_\pi T_\pi
\]
according as \(A=B\), \(|A\cap B|=1\), or \(A\cap B=\varnothing\).

\begin{proposition}[Exact rooted factorization]\label{prop:factorization}
The doubled part of the fixed-five boundary signature
\(\mathcal C_5(P)\) is
\[
\begin{split}
\mathcal C_5(P)\cap\{AA,AT_\pi,T_\pi T_\pi\}
 =\{&AA:\mathsf E\in\operatorname{rel}(R,S)\}\\
 {}\cup\{&AT_\pi:\mathsf I\in\operatorname{rel}(R,S)\}\\
 {}\cup\{&T_\pi T_\pi:
           \mathsf D\in\operatorname{rel}(R,S)\}.
\end{split}                                                \tag{6.7}
\]
\end{proposition}

\begin{proof}
Given a labelling of \(P\) with doubled word \((A,A,B,B)\), restore
\(e\) with label \(A\) and \(f\) with label \(B\).  Cutting the
three-edge cut gives two rooted labellings with a common ordered
connector triangle.  One global coordinate permutation normalizes it
to \(\tau\), placing the transformed root labels in \(R,S\); equality,
intersection size, and disjointness are invariant under that
permutation.

Conversely, take defining labellings for any \(A\in R\) and \(B\in S\).
Their connector words are identically \(\tau\), so corresponding
semiedges glue.  Delete the two roots while retaining their labels on
the four boundary semiedges.  The resulting word is
\((A,A,B,B)\), with the type determined by the cross-relation.  Both
inclusions follow.
\end{proof}

\begin{lemma}[Sparse cross-relations]\label{lem:sparse-relations}
Let \(R,S\) be nonempty and invariant under the stabilizer of \(\tau\).
\begin{enumerate}[label=(\alph*),leftmargin=2.4em]
\item
\(\operatorname{rel}(R,S)=\{\mathsf E\}\) if and only if
\[
 R=S=\{A\}\quad\text{for some }
 A\in\{01,02,12,34\}.                                  \tag{6.8}
 \label{eq:singleton-root}
\]
\item
\(\operatorname{rel}(R,S)=\{\mathsf D\}\) if and only if, after
possibly swapping \(R,S\),
\[
 \varnothing\neq R\subseteq\{01,02,12\},\qquad
 \varnothing\neq S\subseteq
 \{d\in D_5:d\cap a=\varnothing\text{ for every }a\in R\},
                                                               \tag{6.9}
 \label{eq:disjoint-pattern}
\]
where \(S\) is stabilizer-invariant.  There are thirteen unordered
patterns in \eqref{eq:disjoint-pattern}.
\end{enumerate}
\end{lemma}

\begin{proof}
If every cross-pair is equal, choosing one member on either side forces
both sets to be the same singleton.  The singleton is
stabilizer-invariant exactly for the four fixed labels in
\eqref{eq:singleton-root}, proving (a).

For (b), every pair described in \eqref{eq:disjoint-pattern} is
disjoint.
Conversely, suppose every cross-pair is disjoint.  If one side contains
\(34\), the other side consists of triangle edges.  If one side
contains a mixed edge \(u3\), invariance also puts \(u4\) there, and
the only edge disjoint from both is the triangle edge on
\(\{0,1,2\}\setminus\{u\}\).  Thus, after swapping, one side is a
nonempty subset of \(\{01,02,12\}\), and the other is exactly restricted
as in \eqref{eq:disjoint-pattern}.  For a one-edge triangle subset the common
disjoint neighbourhood has two stabilizer orbits, giving three nonempty
invariant choices; this contributes \(3\cdot3=9\).  A two-edge subset
has neighbourhood \(\{34\}\), contributing three, and the three-edge
subset contributes one.  Hence \(9+3+1=13\).
\end{proof}

\begin{lemma}[Cap roots are nonbridges]\label{lem:cap-nonbridge}
In the cap geometry preceding Proposition~\ref{prop:factorization},
neither \(e\) nor \(f\) is a bridge of its rooted shore.
\end{lemma}

\begin{proof}
Suppose \(e\) is a bridge of \(Q_1\), splitting its proper graph into
connected vertex sets \(X,Y\).  Let \(k\) of the three cut connectors
be incident with \(X\).  If \(k=0\) or \(3\), then \(e\) is a bridge of
\(G\), contrary to bridgelessness.  Thus \(k\in\{1,2\}\).  Choose
\(Z=X\) for \(k=1\), and \(Z=Y\) for \(k=2\).  Exactly two edges of
\(G\) leave \(Z\): \(e\) and one connector.

If \(Z\) has \(n\) vertices and \(m\) internal edges, cubic degree
counting gives \(3n=2m+2\), so
\[
 m-n+1=n/2>0.
\]
Hence \(Z\) contains a cycle.  Its complement is connected: the other
bridge component of \(Q_1\) connects through the two remaining
connectors to the connected opposite shore \(Q_2\).  Applying the same
degree count to the complement gives a cycle there as well.  The two
leaving edges form a cyclic two-edge cut of \(G\), a contradiction.
The argument for \(f\) is symmetric.
\end{proof}

M\'a\v cajov\'a, Mazzuoccolo, and Tabarelli
\cite{MacajovaEtAl2026} isolate the following two exceptional type sets
in their unbounded-colour signature relation:
\begin{equation}
\begin{aligned}
\mathcal E_4(i;j,k)&=\{AA,AT_j,AT_k,T_jT_k\},\\
\mathcal E_5(i;j,k)&=\{T_iT_i,T_jT_j,T_kT_k,T_iT_j,T_iT_k\},
 \qquad\{i,j,k\}=\{2,3,4\}.
\end{aligned}
\tag{6.10}
\end{equation}
Below, the statements ``\(\mathcal C_5(P)=\mathcal E_4\)'' and
``\(\mathcal C_5(P)=\mathcal E_5\)'' use the same type masks as
conditional targets in the fixed-five model.  They do not identify
\(\mathcal C_5(P)\) with the full published relation.
Intersecting these displayed sets with the three types in
Proposition~\ref{prop:factorization} gives the following literal
consequence.

\begin{corollary}[Exceptional rooted obligations]
\label{cor:exceptional-rooted}
\begin{enumerate}[label=(\roman*),leftmargin=2.4em]
\item If \(\mathcal C_5(P)=\mathcal E_4(i;j,k)\) and
\(\pi=i\), then
\(\operatorname{rel}(R,S)=\{\mathsf E\}\).  Thus
\(R=S=\{A\}\) for one fixed label
\(A\in\{01,02,12,34\}\), both roots are nonbridges, and every cycle
through either root has a support classified by
Theorem~\ref{thm:translation-obstruction}.
\item If \(\mathcal C_5(P)=\mathcal E_4(i;j,k)\) and
\(\pi\in\{j,k\}\), then
\(\operatorname{rel}(R,S)=\{\mathsf E,\mathsf I\}\).
\item If \(\mathcal C_5(P)=\mathcal E_5(i;j,k)\), then the
cross-relations are disjointness only:
\[
 \operatorname{rel}(R,S)=\{\mathsf D\}.
\]
Hence one of the thirteen patterns in
Lemma~\ref{lem:sparse-relations}(b) occurs.
\end{enumerate}
\end{corollary}

\begin{proof}
For \(\mathcal E_4\), the intersection with
\(\{AA,AT_\pi,T_\pi T_\pi\}\) is \(\{AA\}\) when \(\pi=i\), and
\(\{AA,AT_\pi\}\) when \(\pi=j\) or \(k\).  For
\(\mathcal E_5\), it is \(\{T_\pi T_\pi\}\) for every \(\pi\).
Proposition~\ref{prop:factorization}, Lemmas
\ref{lem:sparse-relations} and \ref{lem:cap-nonbridge}, and
Theorem~\ref{thm:translation-obstruction} give the assertions.
\end{proof}

\subsection{Missing base pairs force proper-core cuts}

For a normalized labelling \(\phi\) and a coordinate \(a\), let
\[
 E_{\bar a}(\phi)=
 \{e\text{ a proper edge of }Q:a\notin\phi(e)\},
 \qquad
 h_a=\{0,1,2,3,4\}\setminus\{a\}.
\]

\begin{lemma}[Complement-cycle switch]\label{lem:complement-switch}
Suppose \(a\notin\phi(r)\).  If \(r\) lies on a proper cycle
\(C\subseteq E_{\bar a}(\phi)\), then replacing every label
\(\phi(e)\), \(e\in C\), by \(\phi(e)\mathbin\triangle h_a\) gives
another normalized labelling.  Consequently, if
\(\phi(r)\mathbin\triangle h_a\notin\mathcal R_5(Q,r)\), then \(r\)
is a bridge of \(E_{\bar a}(\phi)\).
\end{lemma}

\begin{proof}
Every label on \(C\) is a two-subset of the four-set \(h_a\);
symmetric difference with \(h_a\) is its complementary two-subset and
therefore remains in \(D_5\).  At a cycle vertex the same vector
\(h_a\) is added twice, so the vertex equation is unchanged.  No
boundary semiedge is changed.  The last assertion uses the elementary
fact that an edge of a finite graph is a nonbridge exactly when it lies
on a cycle.
\end{proof}

Define the three base pairs
\[
 \mathcal P_0=\{12,03,04\},\qquad
 \mathcal P_1=\{02,13,14\},\qquad
 \mathcal P_2=\{01,23,24\}.
\]

\begin{theorem}[Base-pair cut certificate]
\label{thm:base-pair-certificate}
Suppose \(\mathcal R_5(Q,r)\) is nonempty, contains a label other than
\(34\), and contains none of
\(\mathcal P_0,\mathcal P_1,\mathcal P_2\) as a subset.  Then some
realizing labelling \(\phi\) and \(a\in\{3,4\}\) make \(r\) a bridge of
\(E_{\bar a}(\phi)\).

More precisely, an inner root label \(12,02,\) or \(01\) forces this
bridge for both \(a=3\) and \(a=4\); an outer orbit without its inner
mate supplies one bridge for each coordinate after applying the
symmetry \(3\leftrightarrow4\).
\end{theorem}

\begin{proof}
The fixed boundary word \((01,02,12)\) is invariant under
\(3\leftrightarrow4\), so the exact root signature is invariant as
well.  Complementation in the four-set avoiding \(4\) pairs
\[
 12\leftrightarrow03,\qquad
 02\leftrightarrow13,\qquad
 01\leftrightarrow23,
\]
and complementation in the four-set avoiding \(3\) gives the same
three pairs with \(3,4\) interchanged.  If an inner label occurs,
failure of its base pair and signature invariance exclude both outer
mates.  If an outer orbit occurs, failure of its base pair excludes
the inner mate.  Lemma~\ref{lem:complement-switch} gives the claimed
bridges.
\end{proof}

Here is the promised ordinary cut certificate.  Let
\(\delta_Q^{\mathrm p}(X)\) denote the edge cut in the proper core of
\(Q\), explicitly excluding its three boundary semiedges.  If \(X\)
is one component of \(E_{\bar a}(\phi)-r\), then
\[
 \delta_Q^{\mathrm p}(X)=\{r\}\mathbin{\dot\cup}F_a,
 \qquad F_a\subseteq\{e:a\in\phi(e)\}.
\]
For \(a=3,4\), no boundary label uses \(a\).  The coordinate-\(a\)
support on proper edges is therefore Eulerian, so \(|F_a|\) is even.
If \(r\) is a nonbridge of the unlabelled proper core, then
\(F_a\ne\varnothing\).  Thus Theorem
\ref{thm:base-pair-certificate} supplies such a nonbridge root with an
odd proper-core cut of size at least three, containing \(r\), whose
other edges carry one common outer coordinate.

For completeness, if \(\mathcal R_5(Q,r)=\{34\}\), complementation in
the four-sets avoiding \(0,1,2\) shows that \(r\) is simultaneously a
bridge of \(E_{\bar0}(\phi)\), \(E_{\bar1}(\phi)\), and
\(E_{\bar2}(\phi)\).  Boundary semiedges use these coordinates, so the
even-\(|F_a|\) conclusion is not asserted in this case.

\begin{remark}[Exact limitation]
Corollary~\ref{cor:exceptional-rooted} neither realizes nor excludes any
of the remaining root signatures.  In particular,
Theorem~\ref{thm:translation-obstruction} does not prove that a
nonbridge root has two possible labels: it only classifies what every
root cycle must look like if the label remains unique.  The cut
certificate above is likewise necessary, not contradictory.  The
\(\{\mathsf E,\mathsf I\}\) and thirteen disjoint cases receive no
further elimination here.
\end{remark}

\section{Fixed-five rooted finite frontiers}

\subsection{A human-checkable cap reduction}

Let \(P\) be a connected simple proper core with four specified
terminals \(T=\{t_1,t_2,t_3,t_4\}\), each of proper degree two, and all
other vertices of proper degree three.  Attaching one semiedge at each
terminal gives a terminal-distinct cubic four-pole.  A perfect matching
of \(T\) is \emph{simple for \(P\)} if neither of its pairs is already
an edge of \(P\).

\begin{lemma}[Simple cap matching]\label{lem:simple-cap-matching}
At least one of the three perfect matchings of \(T\) is simple for
\(P\).
\end{lemma}

\begin{proof}
Suppose each perfect matching contains an edge of \(P[T]\), and choose
one such edge from each of the three matching classes of \(K_4\).  The
three chosen edges form either a three-edge star or a triangle.  A star
would give its terminal centre three proper neighbours, although that
terminal has degree two.  In the triangle case, all three triangle
vertices have both proper incident edges inside the triangle, making
it a component of \(P\) disjoint from the fourth terminal.  Both
possibilities contradict the setup.
\end{proof}

The non-Tait consequence used next retains the fixed-five scope.
M\'a\v cajov\'a, Mazzuoccolo, and
Tabarelli~\cite[Lemma 3.8]{MacajovaEtAl2026} prove that a
Tait-colourable pole has no isolated vertex in its signature graph.  The \(T_i\) vertex is
isolated in \(\mathcal E_4\), and \(A\) is isolated in
\(\mathcal E_5\).  Moreover, the cited construction uses at most four
colours, so the same exclusion holds when the exact signature is
restricted to one fixed five-set.

\begin{theorem}[Simple-cap enumeration reduction]
\label{thm:simple-cap-reduction}
Suppose no proper edge of \(P\) is a bridge and
\(\mathcal C_5(P)\) is one of the fixed-five target masks
\(\mathcal E_4,\mathcal E_5\) in \((6.10)\).  Some simple matching on
\(T\) adds two independent edges \(e,f\) to produce a connected simple
bridgeless cubic graph \(G=P+\{e,f\}\) which is not Tait-colourable.
Deleting \(e,f\) recovers \(P\).
\end{theorem}

\begin{proof}
Choose the matching by Lemma~\ref{lem:simple-cap-matching}.  Its pairs
are nonedges with four distinct endpoints, so adding them preserves
simplicity and makes every terminal cubic.  The graph remains
connected.  Every old edge lies on a cycle because the proper core has
no bridge.  Each new edge lies on a cycle because its endpoints are
joined by a path in connected \(P\).  Thus \(G\) is bridgeless.

If \(G\) had a Tait colouring, deleting \(e,f\) and retaining their
colours on the four resulting semiedges would Tait-colour \(P\),
contrary to the fixed-five consequence of the cited lemma.
\end{proof}

\subsection{Root base pairs through order 17}

A rooted three-pole core of order \(n\) here is a connected simple
graph with exactly three degree-two terminals, all other vertices
cubic, and one distinguished proper edge.  Its three boundary
semiedges are ordered and normalized to \((01,02,12)\).

\begin{theorem}[Finite rooted base-pair theorem]
\label{thm:order17-base-pair}
For every odd \(n\leq17\), every nonempty fixed-five root signature of
a nonbridge root in such an order-\(n\) core contains at least one of
\(\mathcal P_0,\mathcal P_1,\mathcal P_2\).
\end{theorem}

This is a computer-assisted finite theorem.  Canonical input was
generated by
\[
 \texttt{geng -cq -d2 -D3 \(n\) \(m\):\(m\)},\qquad
 m=(3n-3)/2,
\]
which forces exactly three degree-two vertices.  At order 17, eight
canonical shards contain \(654\,676\) cores and \(15\,645\,623\)
nonbridge roots.  Two separately written engines were used: a direct
finite-domain propagator and an incremental SAT encoding with a
one-hot ten-label domain and the complete local support table.  Each
reports
\[
\begin{array}{lr}
\text{empty signatures}&337\,059,\\
\text{signatures containing a base pair}&15\,308\,564,\\
\text{violations}&0,\\
\text{solver calls}&52\,216\,251.
\end{array}
\]
Their deterministic transcript digests agree separately on all eight
shards.  The retained verifier checks generator status, row accounting,
both implementation summaries, and every digest equality.  Lower odd
orders are included in the archived report.

Theorem~\ref{thm:order17-base-pair} is not an induction and supplies no
universal root-signature theorem.  Combined with Theorem
\ref{thm:base-pair-certificate}, it says that the ordinary cut
certificate does not occur among nonempty nonbridge roots in this
finite simple-core range because a base pair is always present.

\subsection{The cyclically four-connected order-22 cap corpus}

\begin{theorem}[Finite full-signature cap census]
\label{thm:cyclic4-cap-census}
Let \(G\) be a cyclically \(4\)-edge-connected, bridgeless, non-Tait,
connected simple cubic graph on 22 vertices, and let \(e,f\) be
independent edges.  Then the deletion pole \(G-\{e,f\}\) has all ten
parity-permitted boundary types in its fixed-five signature:
\[
                         \mathcal C_5(G-\{e,f\})
                         \text{ has mask }\texttt{0x3ff}.
\]
\end{theorem}

The quantifier is finite and exhaustive.  Of \(7\,319\,447\) canonical
connected simple cubic graphs on 22 vertices, \(7\,187\,627\) are
bridgeless.  Independent direct edge-colouring and perfect-matching
filters agree on exactly \(12\,892\) non-Tait graphs and on the hard
corpus digest
\[
\texttt{230f2cd88e72011d73a40c5f3d7d5f9fb0fca6110de2ec2541581bc50982037c}.
\]
Exact cut classification leaves 31 cyclically \(4\)-edge-connected
graphs.  Deleting every independent edge pair gives \(14\,322\)
ordered pole records.

Two classifiers ask ten questions for each record.  One uses CaDiCaL;
the other is independently written finite-domain code.  Both return
\[
\begin{aligned}
 \text{rows}&=14\,322,\\
 \text{queries}&=\text{SAT answers}=143\,220,\\
 \text{nonfull rows}&=0.
\end{aligned}
\]
Their uncompressed output streams are byte-identical, with SHA-256
\[
\texttt{52edd8a189012acefd3565a79d5c58b6f209bbe40156560eb57e2c569fd7c091}.
\]
These checks prove Theorem~\ref{thm:cyclic4-cap-census} conditional on
the generator and two implementations.

\begin{corollary}[Exact finite cap consequence]
No bridge-free connected simple terminal-distinct order-22 four-pole
with fixed-five signature \(\mathcal E_4\) or \(\mathcal E_5\) can
have a cyclically \(4\)-edge-connected simple two-edge cap.
\end{corollary}

\begin{proof}
Theorem~\ref{thm:simple-cap-reduction} places such a cap in the
bridgeless non-Tait order-22 corpus.  Theorem
\ref{thm:cyclic4-cap-census} says the recovered deletion pole has the
full ten-type signature, not either exceptional proper subset.
\end{proof}

This corollary does not cover a bridge-bearing or repeated-terminal
pole, a nonsimple core, a cap with a cyclic two- or three-cut, higher
orders, or the published arbitrary-colour signature relation.  A
larger all-cap order-22 run was still in progress when this draft was
frozen and is not claimed here.

\section{The order-22 computational census}

\subsection{The exact property}

For a Tait-colourable cubic graph \(H\), define a conflict graph
\(\Omega(H)\) whose vertices are the edge objects of \(H\).  Join two
edge objects when they are adjacent in \(H\), or when some Tait
colouring places them on the same bichromatic factor cycle.

\begin{lemma}[Conflict-graph equivalence]\label{lem:conflict}
A four-edge set \(S\) is a universally separated matching if and only if
it is an independent set of size four in \(\Omega(H)\).
\end{lemma}

\begin{proof}
Adjacency conflicts are exactly the negation of the matching condition.
Every other conflict records one colouring and one bichromatic cycle
containing two selected edges.  Absence of all such conflicts is exactly
separation in every Tait colouring.
\end{proof}

\begin{theorem}[Order-22 census]\label{thm:census}
No Tait-colourable connected simple cubic graph on 22 vertices has a
universally separated four-edge matching.
\end{theorem}

\subsection{Canonical generation}

The program \texttt{geng} from nauty and Traces 2.9.3
\cite{McKayPiperno2014} was run in four modular shards:
\begin{verbatim}
geng -cq -d3 -D3 22 0/4
geng -cq -d3 -D3 22 1/4
geng -cq -d3 -D3 22 2/4
geng -cq -d3 -D3 22 3/4
\end{verbatim}
Here \texttt{-c} requires connectedness, \texttt{-d3 -D3} requires
degree three, and \texttt{geng} canonically generates simple graphs.
The four disjoint shards contain respectively
\[
 1\,376\,411,\quad 2\,078\,782,\quad
 1\,716\,645,\quad 2\,147\,609
\]
graphs, for a total of \(7\,319\,447\).

\subsection{Two enumeration methods}

The primary program fixes the three colours at one vertex, thereby
quotienting the six global colour permutations.  It recursively
enumerates every remaining Tait colouring, records each bichromatic
factor cycle, builds \(\Omega(H)\), and performs exact backtracking for
an independent four-set.

The replay program does not use the primary colouring recursion.  It
enumerates every perfect matching \(M\), retains \(M\) exactly when its
complement is a union of even cycles, and chooses independently one of
the two alternating colourings on every complementary cycle.  It then
records the three bichromatic factor partitions and builds the same
semantic conflict graph.  Each Tait colouring modulo global colour
permutation is visited six times: three choices of distinguished colour
matching and two orders for the other two colours.

Both methods skip a graph only after establishing that it has no Tait
colouring.  On every shard they agree on the graph count, the number of
Tait-colourable graphs, the number of normalized Tait colourings, and
the zero witness count:

\begin{center}
\begin{tabular}{@{}rrrrr@{}}
\toprule
Shard & Graphs & Tait-colourable & Normalized colourings & Witnesses\\
\midrule
0 & 1,376,411 & 1,338,944 & 17,481,016 & 0\\
1 & 2,078,782 & 2,029,518 & 28,040,733 & 0\\
2 & 1,716,645 & 1,686,315 & 22,708,101 & 0\\
3 & 2,147,609 & 2,119,958 & 27,130,262 & 0\\
\midrule
Total & 7,319,447 & 7,174,735 & 95,360,112 & 0\\
\bottomrule
\end{tabular}
\end{center}

The replay additionally visits \(312\,583\,931\) perfect matchings.
By Lemma~\ref{lem:conflict}, the exhaustive zero-witness output on all
\(7\,174\,735\) eligible graphs proves Theorem~\ref{thm:census},
conditional on the correctness of the generator and implementations.

\subsection{Controls and limitations}

The two implementations agree on complete lower-order corpora:
\begin{center}
\begin{tabular}{@{}rrrrrr@{}}
\toprule
Order & Graphs & Tait-colourable & Colourings & Perfect matchings
  & Witnesses\\
\midrule
12 & 85 & 80 & 307 & 902 & 0\\
16 & 4,060 & 3,848 & 23,420 & 73,929 & 0\\
18 & 41,301 & 39,687 & 307,606 & 995,256 & 0\\
\bottomrule
\end{tabular}
\end{center}
These three controls were rerun from the present source during the
manuscript audit.

A positive target-three control at order 24 has graph6 encoding
\begin{verbatim}
W`??A???A_@_A_cO?S_Gc`_?@O@@?@OO??_????_??H_A?E
\end{verbatim}
Both programs enumerate 36 normalized Tait colourings, return the same
first separated triple (edge identifiers \(0,1,34\)), and reject a
separated four-set.  This exercises the positive witness path, although
it is not a positive four-mark control.

The computation does not address multigraphs, orders 24 and above, or
the rooted avoiding-cycle property.  The two programs share the
mathematical conflict-graph specification and the graph6 input format;
their enumeration engines are different, but they are not formal
proof-producing verifiers.  Reproduction of the full census is the
available check of the finite claim.

\section{Reproducibility}

The source distribution containing this manuscript stores the two
programs at
\begin{verbatim}
scratch/tait_all_coloring_mark_separation.cpp
scratch/verify_order22_separation_via_matchings.cpp
\end{verbatim}
and the machine-readable run record at
\begin{verbatim}
scratch/order22-universal-four-separation-result.json
\end{verbatim}
Paths are relative to the project root, one directory above this
preprint directory.

The finite truth-table checks supporting the rooted algebra, though not
needed by its displayed proofs, are stored at
\begin{verbatim}
scratch/rooted_cycle_translation_blockers.py
scratch/verify_rooted_cycle_translation_blockers.mjs
scratch/rooted-cycle-translation-blockers-result.json
scratch/rooted_three_pole_signature_factorization.py
scratch/verify_rooted_three_pole_signature_factorization.mjs
scratch/rooted-three-pole-signature-factorization-result.json
\end{verbatim}
The first replay checks all \(2^{10}-1\) nonempty \(K_5\) edge
supports.  The second enumerates all nonempty pairs of stabilizer-orbit
unions and confirms the four equality-only, 26 ordered disjoint-only,
and thirteen unordered disjoint-only cases.

The rooted order-17 package and the completed cyclically-four
order-22 cap slice are stored at
\begin{verbatim}
search/rooted-three-pole-frontier-20260727/
  order17-base-pair-report.json
  artifacts/order17-base-pair/
  verify_order17_base_pair_run.py
search/four-pole-order22-cap-20260727/
  artifacts/order22-cyclic4.g6
  artifacts/order22-cyclic4-poles.g6.gz
  artifacts/cyclic4-full-cadical.tsv.gz
  artifacts/cyclic4-full-csp.tsv.gz
  verify_completed_run.py
\end{verbatim}
For the latter package, the verifier was run only through its
cyclically-four extraction and full-signature checks.  Its subsequent
all-cap phase requires artifacts from the still-running larger census
and is not part of this paper.

The semantic build and run commands are:
\begin{verbatim}
clang++ -O3 -std=c++20 \
  scratch/tait_all_coloring_mark_separation.cpp \
  -o /tmp/tait_all_coloring_mark_separation

geng -cq -d3 -D3 22 SHARD/4 |
  /tmp/tait_all_coloring_mark_separation \
    --target 4 --stop-after-first --progress 250000

clang++ -O3 -std=c++20 \
  scratch/verify_order22_separation_via_matchings.cpp \
  -o /tmp/verify_order22_separation_via_matchings

geng -cq -d3 -D3 22 SHARD/4 |
  /tmp/verify_order22_separation_via_matchings \
    --target 4 --stop-after-first --progress 250000
\end{verbatim}

The frozen SHA-256 identifiers are:
\begin{center}
\small
\begin{tabularx}{\textwidth}{@{}lX@{}}
\toprule
Artifact & SHA-256\\
\midrule
Primary source &
\texttt{a00c2f53a3b95a06a7e8767aae2eb2bfa794fc95eeeba18bcac34aea6a12072a}\\
Replay source &
\texttt{104dca4ac143581a9b306172354428a96913f3fe1701c60b2a018510990ecead}\\
\texttt{geng} 2.9.3 binary &
\texttt{ad2f68adf733dbed7cad543841cfa329740596ee5cd136f17a0a17f6e744f5ad}\\
Primary run binary &
\texttt{a63d8048528e3dfa751049fbad6aec3b4909d1f7508f9e56859edbb5af9a6aa0}\\
Replay run binary &
\texttt{3de46d3c343fb3b23679334f1058104488077386f49b4a2541273aa1e5d7d703}\\
Corrected census report &
\texttt{53faa21acfaf75e6bf1d54d6c7d73150a57207067be4f44e20c55be220c322a2}\\
Corrected run JSON &
\texttt{59d8e99cffb573de6ffa8a4df122099db3fd856d8e87803920e9cf69d6946ca8}\\
\bottomrule
\end{tabularx}
\end{center}

The recorded environment was macOS 26.5.2 on arm64 with Apple clang
21.0.0.  Mach-O linker UUIDs can make recompilations byte-distinct, so
the source hashes and semantic counts are the portable anchors; the
binary hashes identify the executables used in the archived runs.
The JSON record includes all shard counts, commands, controls, and scope
warnings.  Its arithmetic totals were independently recomputed during
this manuscript audit.

\section{Relationship among the results and open boundary}

Theorem~\ref{thm:core} is unrooted: its binary cycle may contain a
prescribed edge \(f\).  Theorem~\ref{thm:gate} is rooted but assumes the
marks are separated in one fixed common-colour Tait colouring and leaves
a cyclic four-cut interface.  Neither theorem implies the other.
Proposition~\ref{prop:elliptic-flow} supplies the exact algebraic
language for the later rooted interface.  Theorem
\ref{thm:translation-obstruction} then restricts a different,
exceptional three-pole interface, but only conditionally on a singleton
root signature; Theorem~\ref{thm:base-pair-certificate} turns a wider
failure into an ordinary cut certificate.  Neither is an elimination.
Theorems~\ref{thm:order17-base-pair} and
\ref{thm:cyclic4-cap-census} establish exact finite frontiers in the
fixed-five simple-core model, not universal signature theorems.
Theorem~\ref{thm:census} says only that a particularly strong
four-mark hypothesis cannot occur at order 22 in a Tait-colourable
simple cubic graph.  It is diagnostic evidence, not a substitute for
the missing four-cut argument.

The marked-cycle human-proof frontier left by this paper is:
\begin{quote}
Analyze the rooted four-pole obtained when the forbidden edge belongs
to an independent cyclic four-edge cut, while preserving the inherited
mark-parity and fixed-colouring constraints.
\end{quote}
The distinct fixed-five signature frontier is to prove universal
base-pair closure or reduce the coordinate-cut certificates, and then
to eliminate the surviving cyclically-two/three-cut cap interfaces
beyond the completed finite slice.  Nothing here classifies the
published arbitrary-colour relation.
Even a resolution of that local statement would still require the
remaining global reductions in a Five-Cycle Double Cover proof.  No
claim about the orientable conjecture is made.

\section*{Acknowledgements}

The author thanks the authors of the cited structural and
prescribed-cycle theorems.  The substantial OpenAI Codex contribution
is described in the disclosure on the first page and is repeated here
to avoid ambiguity: the present text and the underlying project are
AI-assisted research artifacts awaiting independent human review.

\bibliographystyle{plain}
\bibliography{references}

\end{document}
